\documentclass[10pt,a4paper]{amsart}
\usepackage[margin=19mm]{geometry}
\usepackage{amsmath,amssymb,mathtools}
\usepackage[scr=rsfs,cal=euler]{mathalfa}
\usepackage{xfrac}
\usepackage[unicode,hidelinks,bookmarks=false]{hyperref}
\newcommand{\ie}{i.e.\ }
\newcommand{\cf}{cf.\ }
\newcommand{\resp}{respectively\ }
\newcommand{\Subsection}{\subsection}
\providecommand{\nobreakdash}{\nobreak\hbox{-}}
\setlength{\emergencystretch}{2em}
\multlinegap0pt
\usepackage{algorithm}\usepackage{algpseudocode}
\usepackage{amssymb}
\usepackage{mathtools}
\usepackage{xcolor}
\usepackage{tikz}
\newtheorem{theorem}{Theorem}
\newcommand{\E}{\mathbb{E}}
\newcommand{\F}{\textsf{False}}
\newcommand{\Gnhalf}{\mathbb{G}(n,1/2)}
\newcommand{\MCS}{\mathsf{LCIS}}
\newcommand{\N}{\mathbb{N}}
\renewcommand{\Pr}{\mathbb{P}}
\newcommand{\T}{\textsf{True}}
\newcommand{\done}{\textsf{AddedPair}}
\newcommand{\eps}{\varepsilon}
\newcommand{\greedy}{\textsc{Greedy} }
\title{Optimal Hardness of Online Algorithms\\ for Large Common Induced Subgraphs}
\author{David Gamarnik, Mikl\'os Z. R\'acz, Gabe Schoenbach}
\date{Source: arXiv:2605.03893v1 (CC BY 4.0)}
\begin{document}
\maketitle

\noindent\emph{Source and licence.} Optimal Hardness of Online Algorithms\\ for Large Common Induced Subgraphs. Version arXiv:2605.03893v1 (CC BY 4.0), distributed by arXiv under the Creative Commons Attribution 4.0 International licence, http://creativecommons.org/licenses/by/4.0/, as recorded in the arXiv licence metadata for this exact version and shown on the abstract page https://arxiv.org/abs/2605.03893v1. Full licence notice: \texttt{licenses/arxiv-2605.03893-CC-BY-4.0.txt}.\par\smallskip

\noindent\textit{Editorial scope.} Counted unit: the article's theorem thm:algs-formal (source0.tex lines 181--186). The article's complete original proof of it is delivered with it (source0.tex lines 567--720, one \texttt{proof} environment of 154 lines, which the article places away from the statement). No mathematics is written, repaired or re-derived: the statement and the proof are the article's own lines, in the article's own notation, and no retained line is substituted or re-flowed.  Retained uncounted as the source-local context the proof needs: lines 526--564.  Nothing was omitted from any retained span, so the package carries no omission record.  No retained line contains a citation, so the wrapper carries no bibliography and no bibliography entry.  No retained line contains a figure environment, an image-inclusion command or drawing code, so no image is delivered and the package declares no asset dependency.  Editorial formatting only, in addition: an AMS \texttt{amsart} preamble that restates the article's own declarations and macros used by the retained lines, namely the environments \texttt{theorem} on separate counters, together with the standard packages those lines need.  The retained environments print in this document as Theorem 1 in order of appearance, whereas the article prints the same items with the numbering of its own full document.  The complete native source of the article as distributed in the arXiv e-print for this exact version is bundled unchanged as \texttt{sources/199826/source0.tex}.
\begin{algorithm}[ht!]
\textbf{Input:} Two graphs $G_{1}$ and $G_{2}$. \\ 
\textbf{Output:} Two induced subgraphs $H_{1} \subseteq G_{1}$ and $H_{2} \subseteq G_{2}$ such that $H_{1} \cong H_{2}$, and a graph isomorphism $\pi : V(H_{1}) \to V(H_{2})$.
\begin{algorithmic}[1]
\Procedure{Greedy}{$G_1, G_2$} 
    \State Let $u_1, \dots, u_n$ be an arbitrary ordering of $V(G_{1})$. 
    \State Let $v_1, \dots, v_n$ be an arbitrary ordering of $V(G_{2})$.
    \State Initialize $S_1:= \emptyset$ and $ S_{2} := \emptyset$. 
    \State Initialize $\pi : S_{1} \to S_{2}$ to be the empty mapping.
    \For{$i \in [n]$}
        \State $\done := \F$;
        \State $j := 1$;
        \While{$\done == \F$ and $j \leq i$}
            \If{$v_{j} \notin S_{2}$ and $\textsc{IdenticalConnections}(u_{i}, v_{j}, S_{1}, S_{2}, \pi, G_{1}, G_{2})$}
                \State $S_{1} \gets S_{1} \cup \{ u_{i} \}$
                \State $S_{2} \gets S_{2} \cup \{ v_{j} \}$
                \State $\pi(u_{i}) := v_{j}$;
                \State $\done \gets \T$;
            \ElsIf{$u_{j} \notin S_{1}$ and $\textsc{IdenticalConnections}(u_{j}, v_{i}, S_{1}, S_{2}, \pi, G_{1}, G_{2})$}
                \State $S_{1} \gets S_{1} \cup \{ u_{j} \}$
                \State $S_{2} \gets S_{2} \cup \{ v_{i} \}$;
                \State $\pi(u_{j}) := v_{i}$;
                \State $\done \gets \T$;
            \Else 
                \State $j \gets j+1$;
            \EndIf
        \EndWhile 
    \EndFor
    \State $H_{1} := G_{1}[S_{1}]$;
    \State $H_{2} := G_{2}[S_{2}]$;
    \State \Return $(H_{1}, H_{2}, \pi)$
\EndProcedure
\State
\Procedure{IdenticalConnections}{$u, v, S_1, S_2, \pi, G_1, G_2$}
    \State \Return $\forall u' \in S_{1}: \quad (u, u') \in E(G_1) \iff (v, \pi(u')) \in E(G_2)$
\EndProcedure
\end{algorithmic}
\caption{Greedy algorithm for finding common induced subgraphs}\label{alg:greedy}
\end{algorithm}

\begin{theorem}[Algorithms, formal]\label{thm:algs-formal}
    Let $(G_1, G_2) \sim \Gnhalf^{\otimes 2}$, define $\greedy$  as in Algorithm~\ref{alg:greedy}, 
    and let $(H_1, H_2, \pi)$ denote the output of $\greedy(G_1, G_2)$. 
    Then $(H_1, H_2)$ is a solution to $\MCS(n)$ and $\pi$ is a graph isomorphism between $H_{1}$ and $H_{2}$. 
    Furthermore, for any $\eps > 0$ and for all $n$ large enough, the size of $(H_1, H_2)$ is at least $(2 - \eps)\log_2 n$ with probability at least $1 - \exp(-n^{\eps / 4})$.
\end{theorem}

\begin{proof}[Proof of Theorem~\ref{thm:algs-formal}]
Throughout the execution of \greedy the induced subgraphs 
$G_{1}[S_{1}]$ and~$G_{2}[S_{2}]$ 
are isomorphic by construction, 
with $\pi : S_{1} \to S_{2}$ being a graph isomorphism. 
Consequently, this remains true for the output of \greedy too. 

We now turn to the lower bound on the size of the output. 
We first introduce some helpful notation. 
Let $S_{1}(t)$ denote the set $S_{1}$ after \greedy processes the vertex pair $(u_{t}, v_{t})$, 
and define $S_{2}(t)$ analogously. 
For $k \in \N$, let $T_{k} := \min \left\{ t : |S_{1}(t)| = k \right\}$ denote the first time when $S_{1}$ has size~$k$. 
For convenience we set $T_{0} := 0$, note that $T_{1} = 1$, and by convention we set $T_{k} = \infty$ if $|S_{1}(n)| < k$. 
To abbreviate expressions that appear throughout the proof, 
set 
\[
\ell := \left\lceil (2-\eps) \log_2 n \right\rceil
\qquad 
\text{and} 
\qquad 
\widetilde{n} := \left\lfloor \frac{n}{2 \log_2 n} \right\rfloor,
\]
and note that $\ell \widetilde{n} \leq n$ (assuming $\eps \geq 1/\log_{2}(n)$).

With the notation above, the size of the output of \greedy is $|S_{1}(n)|$ 
and our goal is to show that $|S_{1}(n)| \geq (2-\eps) \log_2 n$ 
with high probability. 
Observe that $|S_{1}(n)| \geq (2-\eps) \log_2 n$ if and only if $T_{\ell} \leq n$. 
Note also that if $T_{\ell} > n$, 
then there exists $k \in \{1, \ldots, \ell - 1\}$ such that 
$T_{k} \leq k \widetilde{n}$ and $T_{k+1} > (k+1) \widetilde{n}$. Thus by a union bound we have that 
\begin{equation}\label{eq:error_blocks}
\Pr \left( |S_{1}(n)| < (2-\eps) \log_{2}n \right) 
= \Pr( T_{\ell} > n )
\leq \sum_{k=1}^{\ell-1} \Pr \left( T_{k} \leq k \widetilde{n}, T_{k+1} > (k+1) \widetilde{n} \right).
\end{equation}

In the remainder of the proof we fix $k \in \{1, \ldots, \ell - 1\}$
and bound the probability on the right hand side of~\eqref{eq:error_blocks}. 
To do so, we condition on the sets $S_{1}(k \widetilde{n})$ and $S_{2}(k \widetilde{n})$, obtaining that 
\[
\Pr \left( T_{k} \leq k \widetilde{n}, T_{k+1} > (k+1) \widetilde{n} \right) 
= \E \left[ \Pr \left( T_{k} \leq k \widetilde{n}, T_{k+1} > (k+1) \widetilde{n} \, \middle| \, S_{1}(k \widetilde{n}), S_{2}(k \widetilde{n}) \right)  \right].
\]
Note that 
$\Pr \left( T_{k} \leq k \widetilde{n}, T_{k+1} > (k+1) \widetilde{n} \, \middle| \, S_{1}(k \widetilde{n}), S_{2}(k \widetilde{n}) \right) = 0$ 
unless 
$|S_{1}(k \widetilde{n})| = |S_{2}(k \widetilde{n})| = k$. 
Thus in what follows we may and will assume that 
$|S_{1}(k \widetilde{n})| = |S_{2}(k \widetilde{n})| = k$. 

Let 
$U_{k} := \left\{ u_{k \widetilde{n}+1}, \ldots, u_{(k+1)\widetilde{n}} \right\}$ 
and 
$V_{k} := \left\{ v_{k \widetilde{n}+1}, \ldots, v_{(k+1)\widetilde{n}} \right\}$, 
and note that $|U_{k}| = |V_{k}| = \widetilde{n}$. 
For $u \in U_{k}$, 
let $X_{1,k}(u) \in \{0,1\}^{k}$ be the random variable encoding the connectivity of $u$ to the vertices in $S_{1}(k \widetilde{n})$; that is, the coordinates of $X_{1,k}(u)$ are the indicator random variables $\mathbf{1}{\{ (u,u') \in E(G_{1})\}}$ 
for $u' \in S_{1}(k \widetilde{n})$. 
For $v \in V_{k}$, 
define $X_{2,k}(v) \in \{0,1\}^{k}$ analogously, 
and in particular order the coordinates in a way such that 
if the $i$-th coordinate of $X_{1,k}(u)$ corresponds to the indicator random variable $\mathbf{1} {\{ (u,u') \in E(G_{1})\}}$ for some $u' \in S_{1}(k \widetilde{n})$, 
then the $i$-th coordinate of $X_{2,k}(v)$ corresponds to the indicator random variable $\mathbf{1} {\{ (v,\pi(u')) \in E(G_{2})\}}$. 
With these definitions, observe that 
if $T_{k} \leq k \widetilde{n}$ and $T_{k+1} > (k+1) \widetilde{n}$, 
then for every $u \in U_{k}$ and for every $v \in V_{k}$ 
we have that $X_{1,k}(u) \neq X_{2,k}(v)$ 
(since if it were the case that $X_{1,k}(u) = X_{2,k}(v)$ 
for some $u \in U_{k}$ and $v \in V_{k}$, 
then \greedy would add $u$ to $S_{1}$ and $v$ to $S_{2}$, 
which contradicts $T_{k+1} > (k+1) \widetilde{n}$). 
Since $\{ X_{1,k}(u) \}_{u \in U_{k}}$ and $\{ X_{2,k} (v) \}_{v \in V_{k}}$ are all independent of $S_{1}(k \widetilde{n})$ and $S_{2}(k \widetilde{n})$, we thus have that 
\begin{equation}\label{eq:block_gap}
\Pr \left( T_{k} \leq k \widetilde{n}, T_{k+1} > (k+1) \widetilde{n} \, \middle| \, S_{1}(k \widetilde{n}), S_{2}(k \widetilde{n}) \right) 
\leq 
\Pr \left( \forall u \in U_{k} \ \forall v \in V_{k} : X_{1,k}(u) \neq X_{2,k}(v) \right). 
\end{equation}

To bound the probability in the display above, we can first condition on $\{ X_{1,k}(u) \}_{u \in U_{k}}$ and then use that $\{ X_{2,k} (v) \}_{v \in V_{k}}$ are i.i.d.\ uniform to obtain that 
\begin{align}
&\Pr \left( \forall u \in U_{k} \ \forall v \in V_{k} : X_{1,k}(u) \neq X_{2,k}(v) \right) \notag \\
&\qquad \qquad \qquad \qquad = \E \left[ \Pr \left( \forall u \in U_{k} \ \forall v \in V_{k} : X_{1,k}(u) \neq X_{2,k}(v) \, \middle| \, \{ X_{1,k}(u) \}_{u \in U_{k}} \right) \right] \notag \\
&\qquad \qquad \qquad \qquad = \E \left[ \Pr \left( \forall u \in U_{k} : X_{1,k}(u) \neq X_{2,k}(v) \, \middle| \, \{ X_{1,k}(u) \}_{u \in U_{k}} \right)^{\widetilde{n}} \right] \notag \\
&\qquad \qquad \qquad \qquad = \E \left[ \left( 1 - \frac{\left| \left\{ X_{1,k}(u) : u \in U_{k} \right\} \right|}{2^{k}} \right)^{\widetilde{n}} \right], \label{eq:key_prob_calc}
\end{align}
where in the second-to-last line $v$ is an arbitrary vertex in $V_{k}$. 
In the rest of the proof we bound the quantity in~\eqref{eq:key_prob_calc} in two different ways, depending on whether $k$ is small or large. 

First, assume that $k \leq (1-\eps/3) \log_{2} n$. 
Note that the quantity in the expectation in~\eqref{eq:key_prob_calc} is at most $1$, and it is equal to $0$ precisely when 
$\left| \left\{ X_{1,k}(u) : u \in U_{k} \right\} \right| = 2^{k}$. 
Consequently, we can bound~\eqref{eq:key_prob_calc} by the probability that 
$\left| \left\{ X_{1,k}(u) : u \in U_{k} \right\} \right| < 2^{k}$, 
obtaining that 
\begin{align*}
\E \left[ \left( 1 - \frac{\left| \left\{ X_{1,k}(u) : u \in U_{k} \right\} \right|}{2^{k}} \right)^{\widetilde{n}} \right]
&\leq 
\Pr \left( \left| \left\{ X_{1,k}(u) : u \in U_{k} \right\} \right| < 2^{k} \right) \\
&= \Pr \left( \exists x \in \{0,1\}^{k} : \forall u \in U_{k} : X_{1,k}(u) \neq x \right) \\
&\leq 2^{k} \left( 1 - \frac{1}{2^{k}} \right)^{\widetilde{n}} 
\leq 2^{k} \exp \left( - \frac{\widetilde{n}}{2^{k}} \right) 
\leq n \cdot \exp \left( 1 - \frac{n^{\eps/3}}{2 \log_2 n} \right),
\end{align*}
where in the last line we used a union bound, 
that $\{ X_{1,k}(u) \}_{u \in U_{k}}$ are i.i.d.\ uniform, 
the inequality $1-x \leq e^{-x}$, 
and also that $2^{k} \leq n^{1-\eps/3}$. 


Next, assume that 
$(1-\eps/3) \log_{2} n \leq k \leq (2-\eps) \log_{2} n$. 
In this case we argue that with high probability 
$\left| \left\{ X_{1,k}(u) : u \in U_{k} \right\} \right| \geq n^{1-\eps/2}$, 
which allows us to bound the expectation in~\eqref{eq:key_prob_calc}. 
First, we bound the probability that 
$\left| \left\{ X_{1,k}(u) : u \in U_{k} \right\} \right| < n^{1-\eps/2}$
using a simple union bound as follows: 
\begin{multline*}
\Pr \left( \left| \left\{ X_{1,k}(u) : u \in U_{k} \right\} \right| < n^{1-\eps/2} \right) 
\leq 
\Pr \left( \exists S \subseteq \{0,1\}^{k} : |S| = \left\lfloor n^{1-\eps/2} \right\rfloor, \forall u \in U_{k} : X_{1,k}(u) \in S \right) \\
\leq \binom{2^{k}}{\left\lfloor n^{1-\eps/2} \right\rfloor} \left( \frac{\left\lfloor n^{1-\eps/2} \right\rfloor}{2^{k}} \right)^{\widetilde{n}} 
\leq \left( 2^{k} \right)^{n^{1-\eps/2}} \left( n^{-\eps/6} \right)^{\widetilde{n}} 
\leq n^{2n^{1-\eps/2} - \eps \widetilde{n}/6},
\end{multline*}
where we used that 
$n^{1-\eps/3} \leq 2^{k} \leq n^{2}$. We thus have that 
\[
\E \left[ \left( 1 - \frac{\left| \left\{ X_{1,k}(u) : u \in U_{k} \right\} \right|}{2^{k}} \right)^{\widetilde{n}} \right] 
\leq 
\left( 1 - \frac{n^{1-\eps/2}}{2^{k}} \right)^{\widetilde{n}} + n^{2n^{1-\eps/2} - \eps \widetilde{n}/6},
\]
and using the bound $2^{k} \leq n^{2-\eps}$ and the inequality $1-x \leq e^{-x}$ we have that 
\[
\left( 1 - \frac{n^{1-\eps/2}}{2^{k}} \right)^{\widetilde{n}} 
\leq 
\left( 1 - \frac{1}{n^{1-\eps/2}} \right)^{\widetilde{n}} 
\leq \exp \left( - \frac{\widetilde{n}}{n^{1-\eps/2}} \right) 
\leq \exp \left( 1 - \frac{n^{\eps/2}}{2 \log_2 n} \right).
\]


Finally, putting together~\eqref{eq:error_blocks},~\eqref{eq:block_gap},~\eqref{eq:key_prob_calc}, and the bounds on~\eqref{eq:key_prob_calc} discussed above, we have that 
\begin{multline}\label{eq:greedy_final_error_prob}
\Pr \left( |S_{1}(n)| < (2-\eps) \log_{2}(n) \right) \\
\leq (2 \log_2 n) \cdot 
\max \left\{ n \cdot \exp \left( 1 - \frac{n^{\eps/3}}{2 \log_2 n} \right), 
\exp \left( 1 - \frac{n^{\eps/2}}{2 \log_2 n} \right) + n^{2n^{1-\eps/2} - \eps \widetilde{n}/6} \right\}
\end{multline}
which is at most $\exp\left( - n^{\eps/4} \right)$ 
whenever $\eps \geq \frac{13 \log \log n}{\log n}$ 
and $n$ is large enough. 
\end{proof}

\end{document}
